\begin{multicols}{2}
\section{Evaluation}

\subsection{Beschreibung der Testpersonen}

Die folgenden Ergebnisse stammen von 8 Testpersonen, die von Teammitglieder befragt wurden.
Die Testpersonen sinds nicht Teil des Projektteams.

\begin{itemize}
  \item \textbf{TP1:} Männlich, 26 Jahre, technischer Hintergrund
        (Informatikstudium), gelegentliche Nutzung von Google Maps,
        kein Vorwissen über die App.
  \item \textbf{TP2:} Weiblich, 29 Jahre, nicht-technischer Hintergrund
        (Wirtschaftsstudium), nutzt hauptsächlich Google Maps und Wien Karte,
        kein Vorwissen über die App.
  \item \textbf{TP3:} Weiblich, 52 Jahre, technischer Hintergrund (Elektronik), nutzt hauptsächlich Google  Maps, kein Vorwissen über die App. 
  \item \textbf{TP4:} Männlich, 23 Jahre, Student der Volkswirtschaftslehre und
        Politikwissenschaft (Universität Wien), nutzt Google Maps täglich,
        kein Vorwissen über die App.
  \item \textbf{TP5:} Weiblich, 21 Jahre, Studentin der Publizistik- und
        Kommunikationswissenschaft (Universität Wien), nutzt Google Maps
        gelegentlich für öffentliche Verkehrsmittel, kein Vorwissen über die App.
      \item \textbf{TP6:} Männlich, 24 Jahre, nicht-technischer Hintergrund
(Lehramtsstudium Geschichte und Geographie, Universität Wien).
Nutzt Google Maps regelmäßig.
Kein Vorwissen über die App.
\item \textbf{TP7:} Männlich, 28 Jahre, BSc Elektrotechnik (macht gerade den Master, TU Wien). Nutzt Google Maps und ist mit Android vertraut. Kein Vorwissen über die App.
\item \textbf{TP8:} Männlich, 32 Jahre, Diplom-Ingenieur Informatik. Nutzt Google Maps und ist auch mit Android vertraut. Kein Vorwissen zur App.
\end{itemize}
\end{multicols}

\subsection{Ergebnisse der Aufgaben}

\begin{table}[H]
  \centering
  \begin{tabular}{|l|c|c|c|c|}
    \hline
    \textbf{Aufgabe} & \textbf{1. Karte / Pin} & \textbf{2. Download} & \textbf{3. Suche} & \textbf{4. Einstellungen / Teilen} \\
    \hline
    TP1 erledigt & Ja & Ja & Ja & Ja \\
    \hline
    TP1 Schwierigkeit & 2 / 7 & 3 / 7 & 2 / 7 & 4 / 7 \\
    \hline
    TP2 erledigt & Ja & Ja & Ja & Ja \\
    \hline
    TP2 Schwierigkeit & 3 / 7 & 4 / 7 & 3 / 7 & 5 / 7 \\
    \hline
    TP3 erledigt & Ja & Ja & Ja & Ja \\
    \hline
    TP3 Schwierigkeit & 2 / 7 & 1 / 7 & 1 / 7 & 3 / 7 \\
    \hline
    TP4 erledigt & Ja & Ja & Ja & Ja  \\
    \hline
    TP4 Schwierigkeit & 2 / 7 & 3 / 7 & 2 / 7 & 5 / 7 \\
    \hline
    TP5 erledigt & Ja  & Ja & Ja & Ja  \\
    \hline
    TP5 Schwierigkeit & 3 / 7 & 2 / 7 & 4 / 7 & 6 / 7 \\
    \hline
    TP6 erledigt & Ja & Ja & Ja & Ja \\
    \hline
    TP6 Schwierigkeit & 1 / 7 & 3 / 7 & 3 / 7 & 5 / 7 \\
    \hline
      TP7 erledigt & Ja & Ja & Ja & Ja \\
    \hline
      TP7 Schwierigkeit & 1 / 7 & 1 / 7 & 1 / 7 & 4 / 7 \\
    \hline
      TP8 erledigt & Ja & Ja & Ja & Ja \\
    \hline
      TP8 Schwierigkeit & 1 / 7 & 1 / 7 & 1 / 7 & 2 / 7 \\
    \hline
  \end{tabular}
  \caption{Aufgabenabschluss und Schwierigkeitsbewertung (1 = sehr einfach, 7 = sehr schwierig)}
  \label{tab:tasks}
\end{table}

\begin{multicols}{2}

\subsection{SUS-Ergebnisse}

Der SUS-Score wird nach der Standardformel berechnet:
Für Aussagen 1, 3, 5, 7, 9 wird der Wert um 1 reduziert;
für Aussagen 2, 4, 6, 8, 10 wird der Wert von 5 subtrahiert.
Die Summe wird mit 2{,}5 multipliziert, um einen Wert zwischen 0 und 100 zu ergeben.

\begin{itemize}
  \item \textbf{TP1:} SUS-Score = \textbf{72{,}5} –
        entspricht einer \enquote{guten} Usability (Grenzwert: 68).
  \item \textbf{TP2:} SUS-Score = \textbf{62{,}5} –
        entspricht einer \enquote{akzeptablen} Usability mit
        Verbesserungspotenzial.
  \item \textbf{TP3:} SUS-Score = \textbf{80{,}0} -
        entspricht einer \enquote{exzellenten} Usability (Grenzwert: 74).
  \item \textbf{TP4:} SUS-Score = \textbf{77{,}5} –
        entspricht einer \enquote{guten} Usability.
  \item \textbf{TP5:} SUS-Score = \textbf{60{,}0} –
        entspricht einer \enquote{akzeptablen} Usability mit Verbesserungspotenzial.
  \item \textbf{TP6:} SUS-Score = \textbf{70{,}0} -
        entspricht einer \enquote{guten} Usability (Grenzwert 68).
          \item \textbf{TP7:} SUS-Score = \textbf{87{,}5} -
        entspricht einer \enquote{exzellenten} Usability (Grenzwert: 74).
          \item \textbf{TP8:} SUS-Score = \textbf{85{,}0} -
        entspricht einer \enquote{exzellenten} Usability (Grenzwert: 74).
\end{itemize}

\subsection{Beobachtungen und Interviewergebnisse}

\subsubsection*{Testperson 1}

\textbf{Aufgabe 1} wurde problemlos abgeschlossen.
TP1 erkannte die farbkodierten Pins sofort und tippte direkt auf den
richtigen Pin.
\textit{„Die Farben helfen wirklich, man sieht sofort, dass Orange
Kultur bedeutet."}

\textbf{Aufgabe 2} verlief mit einer kurzen Verzögerung:
TP1 suchte zunächst nach einem Download-Button auf der Kartenansicht,
bevor er den Downloads-Tab fand.
Das Download-Icon innerhalb der Liste wurde dann schnell identifiziert.

\textbf{Aufgabe 3} wurde reibungslos abgeschlossen.
Die Suchleiste im Bottom-Sheet wurde ohne Umwege gefunden.

\textbf{Aufgabe 4} bereitete die meisten Schwierigkeiten:
TP1 fand den Dunkelmodus in den Einstellungen schnell,
benötigte aber mehrere Versuche, um die Swipe-Geste für die
Share-Funktion zu entdecken.
\textit{„Das Wischen hätte ich nicht sofort erwartet –
ein kleines Icon oder ein Hinweis wäre hilfreich."}

Im Abschlussinterview nannte TP1 die Kartenansicht als stärkste Funktion
und die fehlende Rückmeldung beim Download als Verbesserungspunkt.

\subsubsection*{Testperson 2}

\textbf{Aufgabe 1} wurde abgeschlossen, jedoch mit Unsicherheit:
TP2 war anfangs von der Anzahl der Pins auf der Karte überfordert
und tippte zunächst auf den falschen Pin.
\textit{„Es gibt so viele Pins – vielleicht könnte man sie filtern?"}

\textbf{Aufgabe 2} verlief mit deutlichem Mehraufwand:
TP2 bemerkte den Downloads-Tab erst nach einer kurzen Suchphase.
Die Tab-Struktur (\enquote{Entdecken} / \enquote{Installiert})
war zunächst unklar.
\textit{„Ich hätte nicht gedacht, dass \enquote{Entdecken} bedeutet,
dass alle Datensätze dort sind."}

\textbf{Aufgabe 3} wurde abgeschlossen, allerdings suchte TP2 zunächst
nach einem Suchfeld oben im Bildschirm.
Das Bottom-Sheet wurde erst nach einem kurzen Moment gefunden.

\textbf{Aufgabe 4} war die schwierigste Aufgabe:
TP2 fand den Dunkelmodus nach kurzem Suchen.
Die Swipe-Geste wurde ohne Hinweis nicht entdeckt;
die Share-Funktion konnte nur mit einem Hinweis des Testleiters
abgeschlossen werden.
\textit{„Ich hätte erwartet, dass es einen Teilen-Button direkt
beim Datensatz gibt."}

Im Abschlussinterview schätzte TP2 besonders die Offline-Funktion
und wünschte sich einen Filter für die Pins nach Kategorie sowie
sichtbarere Aktionen im Downloads-Screen.

\subsubsection*{Testperson 3}

\textbf{Aufgabe 1} wurde mithilfe der Suchfunktion abgeschlossen. Die Testperson war verwirrt, dass keine Suchvorschläge angezeigt wurden und würde gerne mehr Informationen in der Infobox sehen. Auch sollte der zur Infobox gehörende Pin gehighlightet werden.

\textbf{Aufgabe 2} wurde schnell und mühelos von der Testperson abgeschlossen.

\textbf{Aufgabe 3} wurde zuerst versucht mit der Suchbox im Downloadsmenu abzuschließen, da die Testperson noch von der vorherigen Aufgabe im Downloadmenu war, nach einem Hinweis wurde die Suchbox der Karte benutzt. 

\textbf{Aufgabe 4} war auch für Testperson 3 die schwerste Aufgabe. Obwohl der Dunkelmodus leicht zu finden war, war es nicht intuitiv, die Box im Downloadmenu zu swipen. Die Share-Funktion selber konnte danach aber benutzt werden.
Die Testperson erwartete, dass die Teilfunktion nach einem langen Drücken auf die Box auftauchen würde.

Im Abschlussinterview fand TP3 die App praktisch und wünschte sich, dass ein Tooltip die Swipe-Funktion erklären würde, dass die Suchboxen universell funktionieren und dass Suchvorschläge angezeigt werden.

\subsubsection*{Testperson 4}

\textbf{Aufgaben 1--3} verliefen ohne größere Probleme.
Der Pin war nach der Suche nicht hervorgehoben, was TP4 kurz kommentierte,
aber nicht aufhielt.
Der Tab-Name \enquote{Entdecken} war kurz unklar;
nach dem Download fehlte eine Bestätigung, die TP4 durch Wechsel zu
\enquote{Installiert} selbst kompensierte.

\textbf{Aufgabe 4:} Die Swipe-Geste auf der Datensatz-Karte im
Downloads-Tab wurde ohne Hinweis nicht entdeckt.
\textit{„Swipen hab ich mit Teilen nicht assoziiert."}

Im Abschlussinterview wurde die Kartenansicht mit den Datensätzen als
nützlichste Funktion genannt. Vermisst werden ein Kategoriefilter und
ein Link zur Originalquelle.

\subsubsection*{Testperson 5}

\textbf{Aufgabe 1} wurde mit einem Umweg abgeschlossen; bei hoher
Pin-Dichte wurde zunächst der falsche Pin angeklickt.
TP5 merkte an, dass die App Ortskenntnisse voraussetzt.
\textit{„Damit man gezielt suchen kann, muss man eigentlich in Wien
leben und sich auskennen."}

\textbf{Aufgabe 2} verlief ohne Probleme, der Datensatz war direkt sichtbar.

\textbf{Aufgabe 3:} Das Such-Bottom-Sheet von unten war unerwartet;
TP5 fragte den Moderator ob die Aufgabe korrekt abgeschlossen wurde,
da kein Pin hervorgehoben war.

\textbf{Aufgabe 4:} Die Swipe-Geste wurde ohne Hinweis nicht gefunden.
TP5 suchte nach einem sichtbaren Share-Icon.

Im Abschlussinterview gab TP5 an, die App nicht zu nutzen, da sie
keinen Mehrwert für sich sieht.

\subsubsection*{Testperson 6}

\textbf{Aufgabe 1} wurde problemlos abgeschlossen. TP6 erkannte die
farbkodierten Pins sofort und tippte direkt auf den richtigen Pin.
In der Infobox vermisste er einen Link zur Originalquelle.

\textbf{Aufgabe 2} verlief mit kurzer Verzögerung: TP6 war beim
Tab-Namen \enquote{Entdecken} zunächst unsicher. Nach dem Download
fehlte eine sichtbare Rückmeldung; TP6 wechselte selbstständig zu
\enquote{Installiert} und vergewisserte sich dort.

\textbf{Aufgabe 3} wurde abgeschlossen, jedoch fragte TP6 den
Moderator, ob die Aufgabe korrekt abgeschlossen sei, da nach der
Suche kein Pin hervorgehoben wurde.

\textbf{Aufgabe 4} war auch für TP6 die schwerste Aufgabe. Der
Dunkelmodus wurde schnell gefunden; die Swipe-Geste für die
Share-Funktion wurde ohne Hinweis nicht entdeckt. TP6 versuchte
zunächst, durch langes Drücken ein Kontextmenü zu öffnen.
\textit{„Swipen hätte ich mit Teilen nicht assoziiert."}

Im Abschlussinterview nannte TP6 die Kartenansicht als stärkste
Funktion. Gewünscht werden ein Direktlink zur Originalquelle,
ein Kategoriefilter für die Pins sowie ein Tooltip für die
Swipe-Geste.

\subsection*{Testperson 7}

\textbf{Aufgabe 1} - \textbf{Aufgabe 3} für einen Poweruser wie TP7 waren die ersten Aufgaben kein Problem. 

\textbf{Aufgabe 4} hat sich als schwierigste Aufgabe herausgestellt. Die TP hat das Problem mit trial and error gelöst.

\subsection*{Testperson 8}

\textbf{Aufgabe 1} - \textbf{Aufgabe 3} waren für einen Entwickler der auch schon selbst Apps geschrieben hat kein Problem. 

\textbf{Aufgabe 4} war auch für TP8 am schwierigsten, wurde aber auch relativ schnell gelöst, die Bedienelemente waren nicht an der erwarteten Stelle.

Positiv anzumerken ist, dass unser Design intuitiv genug war, dass bei TP7 und TP8 keine Probleme bei Aufgaben aufgetreten sind, die sie erwartungsgemäß trivial lösen können sollten.

\subsection{Erkenntnisse}

\textbf{Stärken:}
\begin{itemize}
  \item Die farbkodierten Pins wurden von TP1 und TP2 positiv
        bewertet und als intuitiv wahrgenommen.
  \item Die Suchfunktion auf der Karte (Aufgabe 3) funktionierte für
        alle Testpersonen gut.
  \item Das Grundkonzept der App – offene Daten auf einer Karte – wurde
        als nützlich und alltagsrelevant eingestuft.
\end{itemize}

\textbf{Schwächen:}
\begin{itemize}
  \item \textbf{Swipe-to-reveal nicht entdeckbar:} Alle Testpersonen fanden die Wischgeste für Lesezeichen, Teilen und Löschen
        nicht ohne Hilfe, bzw. hatten schwierigkeiten damit. Das ist das deutlichste Usability-Problem.
  \item \textbf{Tab-Benennung im Downloads-Screen:} \enquote{Entdecken}
        war für TP2 und TP4 nicht klar genug – sie erwarteten dort kuratierte
        Empfehlungen, nicht alle Datensätze.
  \item \textbf{Pin-Überlagerung:} Bei hoher Pindichte auf der Karte
        waren einzelne Pins schwer zu antippen.
  \item \textbf{Fehlende Filteroptionen:} TP1 und TP2 wünschten
        sich die Möglichkeit, Pins nach Kategorie ein- und auszublenden.
  \item \textbf{Pins nicht gehighlighted:} TP3, TP4 und TP5 waren
        nach einer Suche unsicher, welcher Pin ausgewählt war, weil kein
        visuelles Highlight gesetzt wurde.
  \item \textbf{Keine Suchvorschläge:} TP3 war nicht sicher ob die Suche funktioniert, da es kein Live Feedback gab.
  \item \textbf{Fehlende Download-Bestätigung:} TP4 erhielt
        nach dem Herunterladen eines Datensatzes keine sichtbare Rückmeldung
        und musste selbstständig zu \enquote{Installiert} wechseln, um den
        Erfolg zu überprüfen.
  \item \textbf{Suchleiste oben erwartet:} TP5 suchte die Suchfunktion
        zunächst am oberen Bildschirmrand – analog zur Google-Maps-Konvention –
        und fand das Bottom-Sheet erst nach Zögern.
      \item \textbf{Historie:} TP8 hat ein Feature vorgeschlagen: man sollte markieren können welche Orte man schon besucht hat.
\end{itemize}

\end{multicols}